% Generated draft: requires mathematical and source review.
\documentclass[11pt]{article}
\usepackage[margin=25mm]{geometry}
\usepackage{fontspec}
\setmainfont{Latin Modern Roman}
\usepackage{amsmath,amssymb,mathtools,unicode-math}
\setmathfont{Latin Modern Math}
\usepackage{enumitem,needspace,xurl,hyperref}
\setlength{\parindent}{0pt}
\setlength{\parskip}{5pt}
\setlength{\emergencystretch}{3em}
\newcommand{\N}{\mathbb{N}}
\newcommand{\Z}{\mathbb{Z}}
\newcommand{\Q}{\mathbb{Q}}
\newcommand{\R}{\mathbb{R}}
\newcommand{\C}{\mathbb{C}}
\newcommand{\F}{\mathbb{F}}
\newcommand{\A}{\mathbb{A}}
\newcommand{\PP}{\mathbb{P}}
\newcommand{\E}{\mathbb{E}}
\newcommand{\eps}{\varepsilon}
\newcommand{\dd}{\,\mathrm{d}}
\begin{document}
\section*{UnsolvedMath ID: 20002663}
\textbf{Source problem code:} catalogue entry
\subsection{Definitions and mathematical statement}
\textbf{Source question and reconstruction.} In one variable, perturbing a free stochastic process \(A(t)\) by a free Cauchy variable is said to replace its drift by a conditional expectation and make it smooth. Is there a several-variable analogue? The extraction corrupts the epsilon glyph and omits a clear drift time differential. The aid's conservative notation is
\[
 C_\epsilon=\epsilon C,\qquad
 k_t^\epsilon=E[k_t\mid W^*(A(t)+C_\epsilon)],\qquad\epsilon>0.
\]
Here \(E\) is a trace-preserving conditional expectation, not the scalar trace itself. The intended stochastic model is of the usual form \(\mathrm dA=\mathrm dS+k_t\,\mathrm dt\), with drift allowed to depend on the past; its adaptedness, integrability, driving-noise conventions and target multivariable smoothness are not fully specified. The following aid result is at fixed time under a restrictive free-coordinate hypothesis.

\textbf{Reported one-coordinate posterior.} Work in a finite von Neumann algebra with faithful normal tracial state \(\tau\). Let \(A_i\) be bounded self-adjoint and \(C_i\) a standard Cauchy-distributed self-adjoint affiliated variable, free from \(A_i\). Unbounded affiliated variables are needed because the Cauchy law is not a bounded or square-integrable variable. Put \(B_i=A_i+\epsilon C_i\), \(\mu_i=\operatorname{law}(A_i)\), and
\[
 P_\epsilon(u)=\frac{\epsilon}{\pi(u^2+\epsilon^2)},\qquad
 (T_i f)(y)=
 \frac{\int f(x)P_\epsilon(y-x)\,\mathrm d\mu_i(x)}
      {\int P_\epsilon(y-x)\,\mathrm d\mu_i(x)}.
\]
The denominator is strictly positive. The aid reports
\[
 E_{W^*(B_i)}[f(A_i)]=(T_i f)(B_i).
\]
For the required square-integrable inputs this is an \(L^2\) contraction, preserving the trace. The real function \(T_if\) is locally real analytic; the aid reports quantitative derivative bounds, but supplies no explicit constants or norm-radius formula. Smoothness of these scalar posterior factors is not an automatically defined global noncommutative smoothness theorem.

\textbf{Reported free-product multivariable special case.} Assume the pair algebras \(\mathcal M_i=W^*(A_i,C_i)\) are mutually free, and let
\[
 \mathcal B=W^*(B_1,\ldots,B_m).
\]
For a reduced word \(f_1(A_{i_1})\cdots f_\ell(A_{i_\ell})\) with bounded centered factors \(\tau(f_j(A_{i_j}))=0\) and adjacent indices distinct, the aid reports
\[
 E_{\mathcal B}\left[\prod_{j=1}^\ell f_j(A_{i_j})\right]
 =\prod_{j=1}^\ell(T_{i_j}f_j)(B_{i_j}).
\]
The output factors remain centered. Freeness gives the exact word norm formula and contraction
\[
 \left\|E_{\mathcal B}\left[\prod_j f_j(A_{i_j})\right]\right\|_2
 =\prod_j\|(T_{i_j}f_j)(B_{i_j})\|_2
 \leq\prod_j\|f_j(A_{i_j})\|_2
 =\left\|\prod_jf_j(A_{i_j})\right\|_2.
\]
Finite reduced-word sums extend by the stated \(L^2\) continuity. The bounded-input framework makes products well defined; extension to arbitrary unbounded products requires additional domain/integrability hypotheses. This is a reduced-free-product conditional expectation, not a componentwise multiplicative map on all noncommutative polynomials. General dependent coordinate tuples are not covered.

\textbf{Reported multiplicativity obstruction.} For a von Neumann subalgebra \(\mathcal B\) and bounded self-adjoint \(X\), the identity
\[
 E_{\mathcal B}(X^2)=E_{\mathcal B}(X)^2
\]
forces \(X=E_{\mathcal B}(X)\in\mathcal B\), by faithfulness of the trace and vanishing \(L^2\) conditional variance. Thus a proposed universally multiplicative componentwise smoother is trivial on such hidden coordinates. This does not obstruct the centered alternating-word formula, which does not assert multiplicativity within a single coordinate algebra. The source's general several-variable smoothing question and path-dependent drift regularity remain beyond this supplied special case.
\subsection{Short English statement}
The aid gives a multivariable Cauchy posterior only when the coordinate pairs are freely independent. Conditional expectation acts factorwise on centered alternating words and contracts their \(L^2\) norms. It cannot be promoted to a multiplicative map on all polynomials unless the hidden variable is already observed.
\subsection{Sources}
\begin{itemize}
\item Local UnsolvedMath registry, permanent ID \(20002663\), catalogue entry, \texttt{aim-probability-notes.json}, \texttt{notes[104]}, \emph{Free analysis}, source question 3: \url{https://aimath.org/WWN/freeanalysis/freeanalysis.pdf}.
\item Supplied verification dates the original PDF to August 24, 2006 and reconstructs epsilon using adjacent Poisson-kernel notation. The record is labelled reconstructed\_unverified; this draft has not visited the PDF.
\item Supplied research summary on Cauchy posterior kernels, reduced free products, derivative control and the conditional-variance obstruction. No independent authoritative theorem citation or explicit derivative constants are supplied. Attempt 1: partial\_result, status at run ``open'', rigor ``proved'', low novelty confidence; registry status partially\_solved. The machine-generated aid is not peer reviewed and requires independent verification before citation.
\end{itemize}
\end{document}
